\section{Ladder cuts and a coefficient lower bound}
\label{ladder:section}

Reduced excursions and balancing estimates for inventory words were
developed by Sheffield~\cite[Section~3.3]{Sheffield2016}.
We use these ideas to obtain a complete matching probability uniform in
the ladder height. Counting successfully matched intervals will then
give the coefficient estimate needed for finite-size conditioning.
The piece lengths need not have finite mean; the arguments below
integrate reduced lengths and use stationarity in the piece index.

Give each burger increment $+1$ and each order increment $-1$.
These increments form a simple symmetric random walk.
A left ladder piece, with law $L$, is obtained by sampling letters backwards
until their sum first reaches $+1$, then writing them in chronological order.
A right ladder piece, with law $R$, is obtained by sampling forwards until
the sum first reaches $-1$.
Both pieces are finite almost surely. Their net burger-minus-order counts
are respectively $+1$ and $-1$.

Fix a seam in a two-sided iid word. For $h\ge1$, take the letters on its
left up to the first backwards hit of $+h$, and the letters on its right
up to the first forwards hit of $-h$.
For each fixed $h$, the two resulting words are independent concatenations
of $h$ iid $L$-pieces and $h$ iid $R$-pieces.
Indeed, successive level-hitting times restart the iid sampling; reversing
the order of the $h$ left pieces preserves their joint iid law.
Our quantitative goal is the following bound on the finite-map weights.

\begin{proposition}[Coefficient lower bound]
\label{ladder:lower}
There is $c>0$ such that, for all real $x<1$ sufficiently close to $1$,
\begin{equation}
 \sum_{n\ge1} n a_n u^n x^n\ge c(1-x)^{-1/2}.
 \label{ladder:coefficient-bound}
\end{equation}
\end{proposition}

To prove this bound, we first show that the entire seam interval reduces
to the empty word with probability bounded below uniformly in $h$.
We construct boundaries that give burger-only reductions on the left
and orders-only reductions on the right, then compare their type counts.

\subsection{Stationary cuts}

We index pieces by $\mathbb Z$, with boundary $i$ immediately before piece $i$.
Thus the word between boundaries $a<b$ is the concatenation of pieces
$a,\ldots,b-1$.
We first record the averaging fact needed for observations that can depend
on infinitely many pieces.

\begin{lemma}[Stationary averaging]
\label{ladder:averaging}
Let $(Z_i)_{i\in\mathbb Z}$ be iid, and let $Y_i$ be the translates of an
integrable measurable function of the entire sequence $(Z_i)_{i\in\mathbb Z}$.
Then the averages of $Y_i$ in each index direction converge almost surely
to $\E Y_0$.
In particular, a shift-covariant set of boundaries having positive
probability at boundary $0$ has positive limiting density in both
directions. If $0$ belongs to this set, the last such boundary before $k$
is $k-o(k)$, and the first one after $-k$ is $-k+o(k)$.
\end{lemma}

\begin{proof}
An integrable function of a product sequence can be approximated in $L^1$
by bounded functions of finitely many coordinates.
The translates of a bounded finite-window function split into finitely
many iid subsequences. Their averages converge almost surely: for bounded
centered iid variables, the fourth moment of a partial sum of length $r$
is $O(r^2)$, so the deviation probabilities for their averages are summable.

For completeness, the approximation errors are controlled by the following
maximal inequality for any stationary nonnegative sequence $(U_i)$:
\begin{equation}
 \P\left(\sup_{r\ge1}\frac1r\sum_{i=1}^r U_i>\eta\right)
 \le \frac{\E U_0}{\eta},\qquad \eta>0.
 \label{ladder:maximal}
\end{equation}
To prove it, first allow witnessing intervals of length at most $s$.
Among violating starts in $1,\ldots,m$, select an interval at the first
such start, skip the starts it covers, and repeat.
The selected intervals are disjoint, cover every violating start, and lie
in $1,\ldots,m+s$.
Their sum exceeds $\eta$ times their total length, which is at least the
number of violating starts.
Taking expectations, dividing by $m$, and then letting $m$ and $s$ tend to
infinity proves Equation~\eqref{ladder:maximal}.
Apply this inequality to the absolute errors of the finite-window
approximations. As their $L^1$ errors tend to zero, the almost-sure
convergence of the averages follows. Reversing the indices proves the
other direction.

Apply the result to the indicator of the boundary set.
If $b_k\le k$ is its last boundary before $k$ and $d>0$ is its density,
the absence of boundaries between $b_k$ and $k$ gives
$d(k-b_k)=o(k)+o(b_k)=o(k)$.
The negative-index assertion is identical.
\end{proof}

\begin{definition}
For a two-sided iid sequence of $L$-pieces, a boundary $a$ is an
\emph{$L$-cut} if, for every $b>a$, the word between $a$ and $b$ has no
unfilled orders when evaluated from an empty supply.
For a two-sided iid sequence of $R$-pieces, a boundary $b$ is an
\emph{$R$-cut} if, for every $a<b$, the word between $a$ and $b$ leaves
no surviving burgers.
\end{definition}

\begin{lemma}[Positive probability of cuts]
\label{ladder:cuts}
Each fixed boundary has positive probability of being an $L$-cut and
positive probability of being an $R$-cut, in the respective piece laws.
Consequently both cut sets have positive limiting densities in both
index directions.
\end{lemma}

\begin{proof}
We use the strict drift estimate differently for the two cut probabilities.

\paragraph{Shortening left pieces.}
Recall the backwards stopping time $J$ and the integer $K$ from
Proposition~\ref{word:drift}.
At time $J$, the reduced word consists of one burger and $K$ fixed orders
of the opposite type, and $\E K<1$.
Parse backwards from the end of an $L$-piece using successive independent
copies of this stopping rule.
Before the end of any parsing step its relative sum is nonpositive:
a positive sum would force a surviving burger and hence an earlier stop.
The first hit of $+1$ therefore occurs at a parsing-step endpoint.
The number $N_L$ of steps is the first hit of $1$ by a walk whose iid
increments have law $1-K$.
Writing $\mu=1-\E K>0$, bounded stopping gives
\[
 \mu\,\E(N_L\wedge m)
   =\E\left[\sum_{j=1}^{N_L\wedge m}(1-K_j)\right]\le1.
\]
The bound holds because the increments are at most $1$ and the walk is
stopped on its first hit of $1$.
Thus $\E N_L\le1/\mu<\infty$.

Replace each parsing step by its reduction and put these reductions back
in chronological order.
Lemma~\ref{word:reduction} allows this replacement inside any larger word,
including with an older supply of burgers.
The resulting shortened $L$-piece has only burgers and fixed orders,
and its length $\ell$ satisfies
\begin{equation}
 \E\ell
 =\E\left[\sum_{j=1}^{N_L}(1+K_j)\right]
 =\E N_L\,\E(1+K)<\infty.
 \label{ladder:short-length}
\end{equation}
Here reaching parsing step $j$ depends only on the preceding steps, which
justifies the second equality.
The two net type counts in a shortened piece are integrable, sum to $1$,
and have the same law under exchange of the types.
Each therefore has expectation $1/2$.

Across successive forward $L$-pieces, the net count of each type has
positive drift.
Within a shortened piece, the deviation from its initial count is at most
its length. Moreover, for iid lengths of finite mean, $\ell_i/i\to0$
almost surely, by the integrable-tail bound and Borel--Cantelli.
The running infimum of each type count over all shortened letters is
therefore finite almost surely.
A sufficiently large deterministic starting supply of both types
prevents any unfilled order forever with positive probability.
A finite prefix of prescribed single-burger $L$-pieces creates this supply
with positive probability, independently of the remaining tail.
Starting with that prefix and an empty supply proves positive probability
of an $L$-cut at boundary $0$.
Only the shortened length in Equation~\eqref{ladder:short-length} was
integrated; no moment of the original ladder-piece length is needed.

\paragraph{Reversing right pieces.}
For an $R$-cut, we will show that a backward scan from a piece endpoint
has positive probability of never revealing a surviving burger.
Scan the letters backwards from boundary $0$ in a two-sided iid
$R$-piece sequence. Let $S_j$ be the burger-minus-order sum of the first
$j$ scanned letters. This scan has a different law from ordinary iid
backwards sampling.
The density of its first $m$ letters relative to ordinary iid backwards
sampling is
\begin{equation}
 2(-S_m)\,\1_{\{S_j<0\text{ for }1\le j\le m\}}.
 \label{ladder:reverse-density}
\end{equation}
To see this, use the last $h$ pieces for any $h>m$.
Their concatenation is an iid word stopped at its first hit of $-h$.
Write $W$ for its forward walk, started at $0$, and $\tau$ for this
hitting time. The reversed sums satisfy
\[
 S_j=W_\tau-W_{\tau-j}=-h-W_{\tau-j},\qquad 1\le j\le m.
\]
Since $W$ stays above $-h$ before $\tau$, a possible final string read
backwards has strictly negative partial sums.
If its total is $-r$, its chronological start is at
$-h+r$, where $1\le r\le m<h$.
Before killing at $-h$, the expected number of visits to this level is
$2r$: the walk visits it almost surely, and at each departure the
probability of being killed before returning is $1/(2r)$.
This last probability is one half times the elementary probability
$1/r$ of hitting the lower endpoint before returning from its adjacent
site; a departure upwards returns almost surely.
Each visit followed by the specified admissible string ends exactly at
the killing time. Summing over visits and multiplying by the iid
probability of that string proves Equation~\eqref{ladder:reverse-density}.

\paragraph{Removing the terminal order.}
The first backwards letter is necessarily an order.
After removing it, and restarting the sums at $0$, the remaining scan
has finite-dimensional density
\begin{equation}
 (1-S_j)\,\1_{\{\max_{1\le i\le j}S_i\le0\}}
 \label{ladder:remainder-density}
\end{equation}
relative to iid backwards sampling.
Indeed, the first order contributes $-1$; factoring its conditional order
law out of Equation~\eqref{ladder:reverse-density} leaves exactly
Equation~\eqref{ladder:remainder-density}. Explicitly, the original
probability of a specified order is $\tfrac12$ times its conditional
order probability; the $\tfrac12$ cancels the factor $2$ in the density.
The remaining density has no dependence on the removed order's type.

Apply the stopping rule $J$ to this remaining scan.
Under ordinary sampling, on $\{J=j\}$ the previous sums are nonpositive
and $S_j=1-K$.
The density in Equation~\eqref{ladder:remainder-density} on this event is
$K$: when $K=0$ the sign restriction fails and the density is zero;
when $K\ge1$ it is $1-S_j=K$.
Summing over the disjoint events $\{J=j\}$ shows that the probability,
under the remaining-scan law, of ever revealing a surviving burger is
$\E K$. This is a sum of finite-cylinder identities over disjoint
finite values of $J$, so countable additivity suffices; no optional
stopping of an unbounded density process or finite mean of $J$ is used.
With positive probability $1-\E K$, every finite scanned word instead
leaves no burgers. Appending the removed terminal order cannot create a
burger. In particular, every whole-piece word ending at boundary $0$
leaves no burgers, so this event implies an $R$-cut there.

Finally the cut indicators are shift-covariant functions of iid piece
sequences. Their positive densities follow from
Lemma~\ref{ladder:averaging}.
\end{proof}

\subsection{Stationary outputs and matching}

Between consecutive $L$-cuts $a<b$, the reduction is a burger-only word
of length $b-a$, since each piece has net count $+1$.
Assign these burgers in chronological order to the slots
$a,\ldots,b-1$.
The cut set is unbounded in both directions almost surely, so this
defines a stationary sequence of single burgers on all slots.
Similarly, between consecutive $R$-cuts $a<b$, assign the orders-only
reduction of length $b-a$ to these slots in chronological order.
Both output sequences are measurable shift-covariant functions of their
respective iid piece sequences; the cuts need not be independent.

The matching criterion below identifies the counts we must control:
the oldest available burgers and the orders that must still be filled
at the end.

\begin{lemma}[Bottom supply and terminal demand]
\label{ladder:matching}
Let a burger stack and an order word both have length $h$.
Suppose that, for every $1\le k\le h$ and each type $a$, the number of
fixed type-$a$ orders among the terminal $k$ orders is at most the number
of type-$a$ burgers among the original bottom $k$ burgers.
Then the greedy inventory rule matches all orders and leaves no burgers.
\end{lemma}

\begin{proof}
Suppose a first failure occurs, necessarily at a fixed order of some
type $a$: before a first failure there is one burger per remaining order,
so a flexible order can always be filled.
If no earlier flexible order consumed type $a$, all earlier consumption
of that type was fixed. The hypothesis at $k=h$ then rules out failure.

Otherwise consider the last earlier flexible consumption of type $a$.
Let it remove the burger at position $j$ in the original stack, counting
from the bottom, and let $k$ be the number of burgers immediately after
this removal.
While that burger was present, no fixed type-$a$ order could remove an
older type-$a$ burger, and no flexible order could remove anything older
than it. Hence every original type-$a$ burger below position $j$ is
still present just after this removal.
Because a flexible order takes the freshest burger overall, every
remaining burger is below position $j$; in particular $k\le j-1$.
The remaining type-$a$ supply is therefore at least its count in the
original bottom $k$ burgers.
By hypothesis this suffices for all fixed type-$a$ orders in the remaining
suffix of length $k$.
Until the alleged failure, no further flexible order consumes type $a$,
by our choice of the last such consumption. A failure is impossible.
\end{proof}

We now obtain uniform estimates for these supply and demand counts
from the stationary outputs.

\begin{lemma}[Output frequencies]
\label{ladder:frequencies}
Almost surely, each burger type has frequency $1/2$ in the stationary
$L$-output.
The stationary $R$-output has almost-sure frequencies
$(1-\rho)/2,(1-\rho)/2,\rho$ for the two fixed types and the flexible
orders, respectively, for a deterministic $\rho>0$.

On the event that $0$ is an $L$-cut, let $B_a(h,k)$ count type-$a$ burgers
among the bottom $k$ burgers of the reduction of the $h$ pieces starting
at $0$, where $1\le k\le h$.
Then, almost surely on this event,
\begin{equation}
 \lim_{k\to\infty}\sup_{h\ge k}
 \left|\frac{B_a(h,k)}k-\frac12\right|=0,
 \qquad a=1,2.
 \label{ladder:left-counts}
\end{equation}
On the event that $0$ is an $R$-cut, let $C_a(h,k)$ count fixed orders of
type $a$ among the terminal $k$ residual orders of the $h$ pieces ending
at $0$.
Then, almost surely on this event,
\begin{equation}
 \lim_{k\to\infty}\sup_{h\ge k}
 \left|\frac{C_a(h,k)}k-\frac{1-\rho}{2}\right|=0,
 \qquad a=1,2.
 \label{ladder:right-counts}
\end{equation}
\end{lemma}

\begin{proof}
Type exchange symmetry and Lemma~\ref{ladder:averaging} give the claimed
$L$-frequencies.
The event of an $R$-cut at $0$ depends only on pieces before $0$.
Independently requiring piece $0$ to be the single letter $F$, an event
of probability $p/2>0$, makes boundary $1$ a cut as well and places an
$F$ in output slot $0$.
Consequently
\[
 \rho\ge\P(0\text{ is an }R\text{-cut})\frac p2>0.
\]
Stationary averaging and type exchange symmetry give the remaining
$R$-frequencies.

For Equation~\eqref{ladder:left-counts}, let $b_k$ be the last $L$-cut
at or before $k$, on the event that $0$ is a cut.
It satisfies $k-b_k=o(k)$.
For every $h\ge k$, the bottom $b_k$ surviving burgers are exactly the
stationary output in slots $0,\ldots,b_k-1$:
the pieces after the cut $b_k$ require no older supply and cannot remove
these burgers, by Lemma~\ref{word:reduction}.
Thus the bottom-$k$ count differs from the corresponding stationary
prefix count by at most $k-b_k$, uniformly in $h$.
This proves Equation~\eqref{ladder:left-counts}.

For Equation~\eqref{ladder:right-counts}, take the first $R$-cut $c_k$
at or after $-k$; then $c_k=-k+o(k)$.
For every $h\ge k$, the whole-piece word from $-h$ to $c_k$ leaves no
burgers, by the definition of the cut at $c_k$.
Its residual orders cannot be served by later burgers and do not alter
the reduction of the suffix from $c_k$ to $0$.
The terminal $-c_k$ orders therefore agree exactly with stationary output
in these slots.
The possible discrepancy among the terminal $k$ orders is at most
$k+c_k=o(k)$, independently of $h$.
This proves Equation~\eqref{ladder:right-counts}.
\end{proof}

\subsection{Matching at a seam}

\begin{proposition}[Uniform success at a seam]
\label{ladder:uniform}
There is $\delta>0$, depending only on $p$, such that for every integer
$h\ge1$ the seam interval formed by the backwards hit of $+h$ and the
forwards hit of $-h$ reduces to the empty word with probability at least
$\delta$.
\end{proposition}

\begin{proof}
Choose
\[
 \frac{1-\rho}{2}<\alpha<\frac12.
\]
By Lemmas~\ref{ladder:cuts} and~\ref{ladder:frequencies}, for a sufficiently
large deterministic integer $B$, each of the following events has positive
probability in its respective stationary piece sequence:
\begin{align*}
 \mathcal G_L={}&\{0\text{ is an }L\text{-cut, and }
 B_a(h,k)\ge\alpha k
 \text{ for }a=1,2\text{ and all }h\ge k\ge B\},\\
 \mathcal G_R={}&\{0\text{ is an }R\text{-cut, and }
 C_a(h,k)\le\alpha k
 \text{ for }a=1,2\text{ and all }h\ge k\ge B\}.
\end{align*}
Indeed, increasing $B$ makes each event exhaust its positive-probability
cut event, up to a null set.
Take an integer $D\ge B$ with $\alpha D\ge B$.
The event $\mathcal G_R$ is measurable with respect to the past pieces,
so requiring the next $D$ pieces, indexed $0,\ldots,D-1$, to be single
flexible orders preserves positive probability.

Use independent stationary $L$- and $R$-sequences.
For the left word take $h$ pieces starting at $0$, and for the right word
take $h$ pieces ending at $D$.
Unconditionally these have exactly the independent ladder-piece laws
of the seam interval for each fixed $h$.
The event consisting of $\mathcal G_L$, $\mathcal G_R$, and the prescribed
$D$ flexible pieces has probability
\[
 \delta=\P(\mathcal G_L)\P(\mathcal G_R)(p/2)^D>0,
\]
independent of $h$.

On this event the left word reduces to $h$ burgers.
The right word reduces to $h$ orders: it is all flexible if $h\le D$;
otherwise it is the orders-only reduction of $h-D$ pieces ending at $0$,
followed by $D$ flexible orders.
Consider its terminal $k$ orders, with $k\le h$.
If $k\le D$, their fixed-type counts vanish.
If $k>D$ and $k-D\ge B$, each fixed-type count is at most
$\alpha(k-D)\le\alpha k$ by $\mathcal G_R$.
If $0<k-D<B$, each count is at most $B\le\alpha D\le\alpha k$.
In the last two cases $k\ge B$, so $\mathcal G_L$ supplies at least
$\alpha k$ burgers of each type in the bottom $k$ positions.
Thus every inequality in Lemma~\ref{ladder:matching} holds, including
the small values of $k$, and the interval reduces to the empty word.
\end{proof}

\subsection{Counting intervals through a seam}

\begin{proof}[Proof of Proposition~\ref{ladder:lower}]
Let $I_h$ be the seam interval in Proposition~\ref{ladder:uniform}.
Reflection for the simple symmetric walk bounds the hitting time of a
level $h$ by
\[
 \P(\tau_h>m)\le C_0\frac{h+1}{\sqrt m}.
\]
Consequently the two hitting lengths defining $I_h$, divided by $h^2$,
are tight uniformly over $h\ge1$.
Choose a deterministic $A<\infty$ large enough that
\[
 \P\bigl(I_h\text{ reduces to the empty word and }|I_h|\le Ah^2\bigr)
 \ge\delta/2,\qquad h\ge1.
\]
This follows by subtracting the length-tail probabilities from the success
probability; no independence between length and success is needed.
The intervals $I_h$ are distinct as $h$ varies, since the first-hitting
endpoints move strictly with the level.

There are exactly $2n-1$ deterministic intervals of length $2n$ using
letters on both sides of the seam.
Every empty-reducing word has even length.
By Proposition~\ref{word:identity} and Tonelli's theorem, the expected sum
of $x^{|I|/2}$ over all empty-reducing finite intervals through the seam is
\[
 \sum_{n\ge1}(2n-1)a_nu^n x^n.
\]
For $H=\lfloor(1-x)^{-1/2}\rfloor$, the distinct intervals
$I_1,\ldots,I_H$ therefore give
\begin{align*}
 \sum_{n\ge1}(2n-1)a_nu^n x^n
 &\ge \frac\delta2\sum_{h=1}^H x^{Ah^2/2}\\
 &\ge \frac\delta2 H\,x^{A/(2(1-x))}.
\end{align*}
The last factor is bounded away from zero as $x\uparrow1$, while
$H\asymp(1-x)^{-1/2}$.
Since $2n-1\le2n$, this proves
Equation~\eqref{ladder:coefficient-bound}.
\end{proof}
